\section*{Fixed-root deficit telescope on a physical restriction tree}

Let $f:\{-1,1\}^n\to\{-1,1\}$ be nonconstant and freeze
\[
 R={\mathsf H(f)\over\mathsf I(f)}.
\]
For every physical descendant $g$ define the fixed-root deficit
\[
 D_R(g)=R\mathsf I(g)-\mathsf H(g).
\]
For a relevant coordinate $i$ of $g$, let $g_i^+$ and $g_i^-$ be the
physical children, put
\[
 \Delta_i\mathsf H(g)=\mathsf H(g)
 -{\mathsf H(g_i^+)+\mathsf H(g_i^-)\over2},
\]
and define the signed local slack
\[
 s_R(g,i)=\Delta_i\mathsf H(g)-R\mathsf I_i(g).
\]

\begin{lemma}[edge identity]
For every physical split,
\[
 \boxed{
 {D_R(g_i^+)+D_R(g_i^-)\over2}=D_R(g)+s_R(g,i).}       \tag{1}
\]
\end{lemma}

\begin{proof}
Physical edge counting gives
\[
 {\mathsf I(g_i^+)+\mathsf I(g_i^-)\over2}
 =\mathsf I(g)-\mathsf I_i(g).
\]
Substitution into the left side of (1), together with the definition of
$\Delta_i\mathsf H(g)$, gives
\[
 R(\mathsf I(g)-\mathsf I_i(g))
 -(\mathsf H(g)-\Delta_i\mathsf H(g))
 =D_R(g)+s_R(g,i).
\]
\end{proof}

Let $T$ be any finite physical restriction tree rooted at $f$.  An internal
node $v$ carries a descendant $f_v$ and is split on a physical coordinate
$i(v)$.  Give every node its reach weight $w_v=2^{-\operatorname{depth}(v)}$.
Leaves may occur at different depths and need not be constant.

\begin{theorem}[tree-slack telescope]
For every such tree,
\[
 \boxed{
 \sum_{\ell\in\operatorname{Leaf}(T)}w_\ell D_R(f_\ell)
 =D_R(f)+
 \sum_{v\in\operatorname{Int}(T)}w_v s_R(f_v,i(v)).}    \tag{2}
\]
Since $D_R(f)=0$, this also yields the exact nonreusing decomposition
\[
 \boxed{
 \mathsf H(f)=
 \sum_\ell w_\ell\mathsf H(f_\ell)
 +R\sum_v w_v\mathsf I_{i(v)}(f_v)
 +\sum_v w_v s_R(f_v,i(v)).}                            \tag{3}
\]
Every revealed physical edge and every terminal entropy occurs once.
\end{theorem}

\begin{proof}
Multiply (1) at node $v$ by $w_v$.  The two child deficits then have weights
$w_v/2$, exactly their reach weights.  Summing over the internal nodes cancels
each nonroot, nonleaf deficit once with opposite signs, leaving (2).
Similarly, iterating the physical influence identity gives
\[
 \sum_\ell w_\ell\mathsf I(f_\ell)
 =\mathsf I(f)-\sum_vw_v\mathsf I_{i(v)}(f_v).
\]
Insert this into (2), use $R\mathsf I(f)=\mathsf H(f)$, and rearrange to get
(3).
\end{proof}

\paragraph{Descendant-critical corollary.}
Suppose $f$ has ratio at least that of every nonconstant physical descendant.
Then $D_R(g)\ge0$ for every descendant $g$ (and $D_R(g)=0$ for constants).
Consequently every stopped tree satisfies
\[
 \sum_vw_vs_R(f_v,i(v))
 =\sum_\ell w_\ell D_R(f_\ell)\ge0.                    \tag{4}
\]
At any plateau node $g$ with $D_R(g)=0$, (1) shows
$s_R(g,i)\ge0$ for every relevant coordinate.  Equality holds exactly when
every nonconstant child has ratio $R$.  This recovers the B11 critical-slack
identity at every node lying on the fixed-root ratio plateau.

The individual slacks in (4) need not be nonnegative away from the plateau:
they are increments of the nonnegative potential $D_R$ under conditional
expectation.  If $T$ is continued to constants, (2) gives
\[
 \boxed{\sum_vw_vs_R(f_v,i(v))=0.}                      \tag{5}
\]
Thus any positive root slack must be cancelled by negative slack deeper in
the tree unless every encountered slack vanishes.  This is the exact reason
that a nodewise nonnegative-refund proof is too strong, while a signed global
tree remains viable.

\paragraph{Cross-route substitution and scope.}
Equation (2) substitutes the B11 influence-weighted child deficit into the
same signed-telescope position occupied by the independently reviewed
algebraic core of the B05--B19 one-point-defect terminal refund.  The
quantities are not identified: $s_R$ is a fixed-root entropy--influence
potential increment, whereas the B05--B19 credit is the complete physical
$L+Q$ packet.  What transfers exactly is the accounting mechanism: signed
internal increments plus an explicit terminal potential, with reach weights
and no repeated edge.

This does not prove FEI.  The coefficient $R$ in (3) is the unknown root
ratio, and a complete tree makes (5) a zero-sum identity.  The next unpaid
lemma is to construct, for every descendant-critical tensor-prime source, a
stopping tree and an additional nonnegative Boolean compensator that bounds
the terminal entropy/deficit without inserting $R$ as an assumed universal
constant.  It must survive the one-point-defect negative-credit chain,
addressing, near tensors, parity erasure, and shared-variable interference.
